% !TEX root = ../my-thesis.tex
\chapter{Einleitung}
\label{sec:intro}

\section{Motivation}
\label{sec:intro:motivation}

\section{Zielsetzung und Forschungsfragen}
\label{sec:intro:goal}

Die Arbeit orientiert sich an den folgenden Forschungsfragen:

\begin{itemize}
  \item \textbf{F1:} Wie unterscheiden sich die betrachteten
    Workflow-Management-Systeme bei der Modellierung eines Pipeline-
    und eines Scatter-Gather-Workflows?
  \item \textbf{F2:} Welches Laufzeitverhalten zeigen die Systeme bei
    der Ausführung dieser Workflows, gemessen an Makespan,
    referenzbasiertem Overhead und Koordinationsmetriken?
  \item \textbf{F3:} Wie lassen sich die Systeme auf einem
    Slurm-basierten HPC-System betreiben, und wie setzen sich die
    Laufzeiten unter Clusterbedingungen zusammen?
  \item \textbf{F4:} Für welche Arten von Workflows und
    Anwendungsfällen erscheinen die einzelnen Systeme besonders
    geeignet?
\end{itemize}

\section{Aufbau der Arbeit}
\label{sec:intro:structure}

Die Arbeit ist wie folgt aufgebaut:

\textbf{Kapitel \ref{sec:basics}} \\[0.2em]
legt die begrifflichen Grundlagen: wissenschaftliche Workflows, die
beiden untersuchten Muster, Workflow-Management-Systeme, der
HPC-Kontext mit Slurm sowie die fünf betrachteten Systeme und die
verwandten Arbeiten.

\textbf{Kapitel \ref{sec:methodik}} \\[0.2em]
beschreibt die Methodik: das Evaluationsdesign, die Auswahl der
Systeme, die qualitativen Vergleichsaspekte, die quantitativen
Messgrößen sowie Benchmark-Workflows, Versuchsplan und Grenzen.

\textbf{Kapitel \ref{sec:lokal}} \\[0.2em]
dokumentiert die Modellierung beider Muster in den fünf Systemen und
die lokale Messreihe als kontrollierte Vergleichsbasis.

\textbf{Kapitel \ref{sec:mogon}} \\[0.2em]
untersucht die Systeme auf dem HPC-System MOGON: Einsetzbarkeit und
Betriebsmodelle, den nativen Slurm-Modus mit Zeitzerlegung, die
dedizierte Messung, die AiiDA-Konfigurationsstudie und den Vergleich
von exklusivem und geteiltem Knoten.

\textbf{Kapitel \ref{sec:diskussion}} \\[0.2em]
diskutiert die Ergebnisse entlang der Forschungsfragen, stellt den
Zusammenhang von Architektur und Messwerten her, gleicht die Befunde
mit der Terminologie von Suter et al. ab und behandelt Eignung,
praktische Herausforderungen und Grenzen.

\textbf{Kapitel \ref{sec:conclusion}} \\[0.2em]
fasst die Arbeit zusammen und gibt einen Ausblick.
